%% main.tex — ZK-Circuit-Aware NAS
%% NeurIPS 2026 submission

\documentclass[10pt]{article}

\usepackage[margin=1in]{geometry}

%% --- Core packages ---
\usepackage[utf8]{inputenc}
\usepackage[T1]{fontenc}
%\usepackage{microtype}  % disabled: not available in this TeX install
\usepackage{amsmath, amssymb, amsthm}
\usepackage{mathtools}
\usepackage{bm}
\usepackage{booktabs}
\usepackage{graphicx}
\usepackage{hyperref}
\usepackage{xcolor}
\usepackage{algorithm}
\usepackage{algpseudocode}
\usepackage{natbib}
\usepackage{url}

%% --- Theorem environments ---
\newtheorem{theorem}{Theorem}
\newtheorem{lemma}[theorem]{Lemma}
\newtheorem{proposition}[theorem]{Proposition}
\newtheorem{corollary}[theorem]{Corollary}
\newtheorem{definition}{Definition}
\newtheorem{remark}{Remark}

%% --- Math macros ---
\newcommand{\calA}{\mathcal{A}}
\newcommand{\calC}{\mathcal{C}}
\newcommand{\bbR}{\mathbb{R}}
\newcommand{\bbE}{\mathbb{E}}
\newcommand{\logrows}{\mathrm{logrows}}
\newcommand{\ZKcost}{\mathrm{ZKcost}}
\newcommand{\struct}{\mathrm{struct}}
\newcommand{\FLOPs}{\mathrm{FLOPs}}

%% ============================================================
\begin{document}

\title{ZK-Circuit-Aware NAS: A Cost Model, Benchmark, and Proxy Study\\
       for Provable Inference}

\author{%
  Anonymous Author(s)\\
  Under Review
}

\maketitle

% ============================================================
\begin{abstract}
% ============================================================
Deploying neural networks in zero-knowledge (ZK) proof systems requires
co-optimizing for both accuracy and circuit complexity --- yet existing Neural
Architecture Search (NAS) methods treat these objectives independently.  We
make three contributions.  First, we show that, for fixed quantization and lookup-range settings,
the ZK circuit layout depends only on the architecture and not on the weights, so
a cost model can be evaluated before any training (Theorem~\ref{thm:det}); on
73 architectures the model matches ezkl-calibrated logrows exactly for 49 and within one for 66
(Spearman $\rho=0.80$, 95\% CI $[0.65,0.91]$; $\rho=0.92$--$0.94$ for ReLU, $x^2$ and $x^2{+}x$
families but only $0.28$ for $x^3{+}x$).
We further show that FLOPs does not determine ZK cost: for equal
multiply-accumulate counts the cost ratio between activation types is set by
the per-activation circuit cost and the layer fan-in (Theorem~\ref{thm:sep}).
Second, we introduce \textbf{zkml-bench}, the first NAS benchmark spanning
polynomial activations ($x^2$, $x^3$, $x^2{+}x$), BN-free designs, and ZK
proof time as a fitness metric.  We evaluate six zero-cost proxies against
this benchmark and show: no single proxy works reliably across all activation
families (NEAR $\rho{=}{+}0.828$ on our CIFAR-10 ReLU architectures and $+0.933$ on
ImageNet-16-120, but $\rho{=}-0.089$ to $-0.624$ on $x^2$ --- same proxy,
same code, activation type determines reliability; NASWOT, which scores ReLU
activation patterns, is provably degenerate for $x^2$ (our implementation returns $-\infty$ on every
polynomial network because it finds no ReLU modules, so the sweep does not test the proof); ZenNAS is anti-correlated in every
family, with a scaling argument in Proposition~\ref{prop:zen}), while the structural cost proxy tracks ezkl
logrows at $\rho{=}0.80$.  Third, a width-matched ezkl measurement
gives a negative result for a common assumption: ReLU and $x^2$ networks have nearly identical
circuit sizes (equal logrows; $x^2$ has at most $6\%$ more assigned cells), so activation choice is not
the dominant cost driver at these widths, while $x^3{+}x$ networks calibrate to
$\logrows{=}23$ in all unnormalized skeletons we tried.  Using the cost model, within a budget of at most $2^{21}$ estimated circuit constraints
(structural-proxy $\logrows \le 21$), joint-proxy-guided search selects a
$75.0\%$ CIFAR-10 architecture versus $62.8\%$ for random selection ($+12.2$ points),
using $20\times$ fewer constraints than the most accurate ReLU baseline; note that
ezkl calibration places this particular architecture at $\logrows{=}22$
(Table~\ref{tab:ezkl}).  A targeted validation of
26 architectures under three seeds confirms that 30-epoch rankings predict
independently trained 100-epoch rankings ($\rho=0.932$); the selected
$75.0\%$ architecture reproduces at $74.7\%\mathbin{\pm}1.6\%$.
\end{abstract}

% ============================================================
\section{Introduction}
\label{sec:intro}
% ============================================================

\paragraph{The problem.}
ZK-ML systems (ezkl, Orion, RISC Zero) let anyone prove neural network
inference was correct --- without revealing inputs.  This enables private
biometric verification (e.g., Worldcoin), on-chain AI oracles, and
compliant ML auditing.  But proving is slow: in our measurements, real proving
time ranges from under a minute ($47$--$58$s) to over twenty minutes ($1340$s)
depending on architectural choices that practitioners cannot currently predict.

\paragraph{What we show empirically.}
Real proving time in our measurements ranges from under a minute to over twenty minutes
depending on the architecture (Table~\ref{tab:ezkl}).  A width-matched comparison
(Table~\ref{tab:matched}) shows that this spread does not come from the ReLU-versus-$x^2$
choice: at identical widths the two have the same calibrated $\logrows$ and $x^2$ has
at most $6\%$ more assigned cells.  The spread in our sweep comes from layer type and width, and from
ezkl's range calibration, which raises $\logrows$ for $x^3{+}x$ networks.  We therefore report a
negative result for the assumption that polynomial activations are cheaper to prove in ezkl.
Choosing DWS+ReLU over a standard CNN reduces single-run mock-proof time from $1119$s to
$534$s at a cost of $6.6$ accuracy points ($75.0\%$ vs.\ $81.6\%$).  The most accurate ReLU
architecture we evaluated needs $\logrows{=}25$ ($2^{25}$ rows), and we recorded no
mock-proof result for it.

\paragraph{Why this requires NAS.}
Architecture search for ZK inference needs to jointly optimize accuracy and
circuit cost --- but no NAS method, proxy, or benchmark treats ZK cost as a
fitness metric.  We show that no single standard zero-cost proxy works
reliably across all polynomial activation families: NEAR swings from
$\rho{=}{+}0.83$ (our CIFAR-10 ReLU architectures) to $\rho{=}-0.089$ (our $x^2$
architectures), ZenNAS is anti-correlated in every family (pooled
$\rho{=}-0.55$), and NASWOT, which scores binary ReLU activation patterns, is
degenerate on polynomial activations (provably so for $x^2$).  While ZiCo partially works for $x^3{+}x$ architectures
($\rho{=}{+}0.54$), it fails for $x^2{+}x$ ($\rho{=}-0.04$) --- no proxy
generalizes.

\paragraph{Our contributions.}
\begin{enumerate}
\item \textbf{Theory}: for fixed quantization and lookup-range settings, the
  circuit layout is a function of the architecture alone, so a cost model is
  available at initialization time (Theorem~\ref{thm:det}).  FLOPs does not
  determine ZK cost (Theorem~\ref{thm:sep}).  A scaling argument explains ZenNAS's anti-correlation for polynomial
  activations (Proposition~\ref{prop:zen}).

\item \textbf{Benchmark}: First NAS benchmark covering polynomial activations
  $\times$ BN-free designs $\times$ 3 datasets (CIFAR-10/100/ImageNet-16-120),
  with the cost model compared against ezkl-calibrated logrows.

\item \textbf{Measurements}: a width-matched ezkl comparison (ReLU $\approx$
  $x^2$ in circuit size; $x^3{+}x$ calibrates to $\logrows{=}23$), and a
  proxy-versus-ezkl comparison on $73$ architectures (exact for $49$, within one for $66$).

\item \textbf{Application}: Joint proxy finds $75.0\%$ accuracy
  ($+12.2$ points vs.\ random) within a $2^{21}$ estimated-constraint budget
  (ezkl-calibrated $\logrows{=}22$ for the selected architecture).
\end{enumerate}

\paragraph{Paper organization.}
Section~\ref{sec:background} reviews ZK-SNARK cost and contrasts with MPC.
Section~\ref{sec:theory} presents the theoretical results.
Section~\ref{sec:benchmark} introduces zkml-bench.
Section~\ref{sec:findings} reports width-matched ezkl cost measurements.
Section~\ref{sec:results} presents empirical results.
Section~\ref{sec:related} surveys related work.
Section~\ref{sec:conclusion} concludes.


% ============================================================
\section{Background}
\label{sec:background}
% ============================================================

\subsection{ZK-SNARK Circuit Cost Model}

In Halo2/PLONK over BN254 (as used by ezkl), the circuit is a rectangular
table of $2^{\logrows}$ rows.  Cost is dominated by three components:
\begin{itemize}
  \item \textbf{Linear layers}: each multiply-accumulate contributes one
    circuit row, giving $O(w_{\mathrm{in}} \times w_{\mathrm{out}})$ rows per
    layer.
  \item \textbf{Lookup-based activations}: non-linearities such as ReLU are
    evaluated through lookup tables whose size grows with the activation's integer
    range; the table is shared across the circuit.  At the widths we test, the
    resulting per-element cost is small relative to the multiply-accumulate rows
    (Section~\ref{sec:findings}).
  \item \textbf{Polynomial activations} ($x^2$): a single degree-2 gate
    constraint, costing $O(\mathrm{width})$ rows.
\end{itemize}
For a fixed quantization scale and lookup range, the circuit structure is
determined at compile time from the architecture specification and does not
depend on weight values.  In ezkl the scale and range, and hence $\logrows$,
are chosen by a calibration pass on data, which can raise $\logrows$ (Section~\ref{sec:findings});
our cost model therefore predicts the circuit size up to this calibration effect.

\subsection{Contrast with MPC Cost Model}

CryptoNAS~\citep{ghodsi2020cryptonas} optimizes for MPC (garbled circuits).
In MPC: linear layers cost $O(1)$ via secret sharing (effectively free), while
ReLU requires oblivious transfer (expensive).  In ZK-SNARKs, linear layers are
\emph{not} free.  An architecture with half the ReLU count but double channel
width is $2\times$ cheaper in MPC but identical in ZK cost.  CryptoNAS
results do not transfer to ZK deployment.


% ============================================================
\section{Theory}
\label{sec:theory}
% ============================================================

\subsection{ZK Circuit Cost Determinism}

\begin{theorem}[ZK Circuit Cost Determinism]
\label{thm:det}
Fix a quantization scale and lookup range.  For a fully-connected or
convolutional network with $L$ layers, activation types
$\alpha = (\alpha_1,\ldots,\alpha_L)$, and width profile
$W = (w_0,\ldots,w_L)$, the circuit layout (columns, selectors, lookup and
permutation structure) is a function of $(L,\alpha,W)$ only and is independent
of the weight values.  Consequently the row count admits the parametric form
\begin{equation}
C(L,\alpha,W) \;=\; \sum_{i=1}^{L} w_i w_{i-1}
  \;+\; \sum_{i=1}^{L} w_i\,\kappa(\alpha_i)
  \;+\; \varepsilon(L,\alpha,W),
\end{equation}
where $\kappa(\alpha)$ is the per-element row cost of activation type $\alpha$
(determined by the backend's compilation rule for $\alpha$, independent of
weights) and $\varepsilon$ collects permutation and fixed per-node overhead.
Here $C$ counts assigned cells (ezkl's \texttt{total\_assignments}); the circuit uses
$2^{\lceil\log_2 C\rceil}$ rows, so $\logrows=\lceil\log_2 C\rceil$, up to the
calibration-dependent adjustment of the scale and lookup range (which, in Section~\ref{sec:findings},
does not occur for ReLU, $x^2$ and $x^2{+}x$ networks but does for unnormalized $x^3{+}x$).  $C$ is computable in $O(L\cdot\max(W))$ time
before any training.
\end{theorem}

\begin{proof}
The Halo2 \texttt{Circuit::configure} call that determines all column counts,
selector polynomials, lookup arguments, and permutation structure takes only
the architecture specification and the fixed quantization/range settings as
input.  Weight values only populate advice columns at proving time and do not
affect circuit structure.  Hence the number of rows consumed by each layer is
a function of $(w_{i-1},w_i,\alpha_i)$; linear layers contribute one row per
multiply-accumulate, and each activation contributes $\kappa(\alpha_i)$ rows
per element, where $\kappa$ is a constant of the compilation rule.  The residual
$\varepsilon$ is not given in closed form; we estimate it empirically (a fixed
per-node overhead and a logrow snapping step), and Section~\ref{sec:results}
evaluates the resulting model against ezkl.
\end{proof}

\begin{corollary}[Initialization-Time Computability]
Unlike latency (requires hardware measurement) or accuracy (requires training),
the ZK cost model of Theorem~\ref{thm:det} can be evaluated at NAS time with no
training and no hardware measurement.  Its agreement with ezkl is empirical
(Section~\ref{sec:results}), not exact.
\end{corollary}

\subsection{Complexity Separation}

\begin{theorem}[FLOPs Does Not Determine ZK Cost]
\label{thm:sep}
Consider two networks with identical width profile $W$ (hence identical
multiply-accumulate counts $M=\sum_i w_i w_{i-1}$) that differ only in
activation type, with per-element costs $\kappa_A<\kappa_B$ and with the
overhead $\varepsilon$ neglected.  Write $S=\sum_i w_i$ and
$\bar f = M/S$ for the average fan-in.  Then the cost ratio is
\[
\frac{C_B}{C_A} \;=\; \frac{\bar f+\kappa_B}{\bar f+\kappa_A},
\]
which lies strictly between $1$ and $\kappa_B/\kappa_A$ and tends to
$\kappa_B/\kappa_A$ as $\bar f/\kappa_A\to 0$.  In particular, for any
$\bar f\le\kappa_A$ we have $C_B/C_A\ge \kappa_B/(2\kappa_A)$.  Since the two
networks have equal FLOPs, FLOPs cannot be a constant-factor approximation to
ZK cost: the factor depends on $\kappa_B/\kappa_A$ and on the fan-in.
\end{theorem}

\begin{proof}
The cost ratio follows from dividing both costs $C_X = M + S\kappa_X$ by $S$.
The ratio is a monotonically decreasing function of $\bar f$, equal to
$\kappa_B/\kappa_A$ at $\bar f=0$ and tending to $1$ as $\bar f\to\infty$.  If
$\bar f\le\kappa_A$ the numerator is at least $\kappa_B$ and the denominator at
most $2\kappa_A$.
\end{proof}

\begin{remark}
In an idealized model where a lookup activation pays for a full table of $2^B$
rows per element, $\kappa_B=2^B$ and the ratio can approach $2^B/d$ for a degree-$d$
polynomial.  ezkl does not behave this way, and our width-matched measurements
(Section~\ref{sec:findings}) find $\kappa_{\mathrm{ReLU}}\approx\kappa_{x^2}$ at the widths we
test: $x^2$ networks have at most $6\%$ more assigned cells than ReLU networks and the same $\logrows$.  The
separation that the theorem allows is therefore \emph{not} observed between ReLU and $x^2$ in
ezkl.  It is a statement about the structure of the cost.  One place where cost differs at equal
multiply-accumulate counts is the $x^3{+}x$ family, whose assignment counts match the others but whose
calibrated $\logrows$ is $23$ in every unnormalized skeleton we tried; there the difference arises
from ezkl's range calibration rather than from rows.
\end{remark}

\subsection{Proxy Failure Modes for Polynomial Activations}

\begin{proposition}[NASWOT Degeneracy for $x^2$ Activations]
\label{prop:naswot}
For any network whose activations are $\sigma(t)=t^2$ (more generally $t^d$ with
$d$ even), the NASWOT score $\log|\det(K)|$ computed on the binary patterns
$\mathbf{1}[\sigma(\cdot)>0]$ is $-\infty$ for every architecture.
\end{proposition}

\begin{proof}
NASWOT computes $K_{ij} = \langle h_i, h_j \rangle$ where
$h_i^{(\ell)} = \mathbf{1}[\sigma(W^{(\ell)} x_i) > 0]$ is the binary
activation pattern.  For $\sigma(t) = t^2$: since $t^2 \ge 0$ always, the
indicator $\mathbf{1}[t^2 > 0] = 1$ for all non-zero inputs (almost surely
under continuous random inputs).  Therefore all samples map to the all-ones
binary code: $h_i^{(\ell)} = \mathbf{1}$ for all $i$.  The kernel matrix $K$
becomes rank-1, so $\det(K) = 0$, and $\log|\det(K)| = -\infty$ for every
architecture with even-degree pure-power activations.
\end{proof}

\begin{remark}
The argument does not cover $x^3{+}x$ or $x^2{+}x$, whose sign patterns are not
constant.  Our NASWOT implementation scores the binary outputs of
\texttt{nn.ReLU} modules and returns $-\infty$ when a network contains none, which
is the case for every polynomial architecture in the sweep.  The $-\infty$
counts for all $272$ polynomial networks (including the $79$ $x^2$ ones) therefore reflect that the
implementation has no ReLU patterns to score; they do not test the proposition, which is a
statement of mathematics only, and they are not a proof of degeneracy for $x^3{+}x$ or $x^2{+}x$; we did not evaluate a
generalization of NASWOT that uses the sign patterns of polynomial pre-activations.
\end{remark}

\begin{proposition}[ZenNAS Score Inflation for Polynomial Activations]
\label{prop:zen}
For a network with polynomial activation $\sigma(t) = t^d$ ($d \ge 2$) and
$L$ polynomial activation layers, the Zen score satisfies:
\[
\mathrm{Zen}(f) = \Theta(d^L \cdot \log \sigma_x),
\]
where $\sigma_x$ is the input scale.  This quantity is a function of
architectural scale $(d, L,$ weight norms$)$, not of classification accuracy.
By this scaling argument, ZenNAS cannot distinguish a large, unstable
polynomial network from an accurate small one, which would produce the anti-correlation with
accuracy we observe in BN-free settings.
\end{proposition}

\begin{proof}[Proof sketch]
For $\sigma(t) = t^d$, a network output magnitude scales as
$\mathbb{E}[\|f(x)\|_2] = \Theta(\sigma_x^{d^L})$ (verified numerically:
slope $2.00$ for $d=2$, $L=1$; slope $4.00$ for $d=2$, $L=2$; slope $2.96$
for $d=3$, $L=1$).  Both terms of the Zen score grow as
$\Theta(d^L \cdot \log \sigma_x)$.  The same super-exponential scaling
$d^L$ that inflates the Zen score causes gradient magnitude explosion:
$\|\nabla_{W_1} \mathcal{L}\| \sim d^L \cdot \sigma_x^{d^L - 1}$.
Architectures with the highest Zen scores (large, deep polynomial networks)
have training failure rates of $56$--$57\%$ (Table~\ref{tab:zen}), anchoring
those architectures at near-random accuracy ($10\%$).  Without batch
normalization to re-scale outputs to unit variance, output magnitude is
anti-informative: the architectures ZenNAS ranks highest are the ones least
likely to train successfully.
\end{proof}

\subsection{Theory-Empirical Connections}

Table~\ref{tab:theory} summarises the four theoretical predictions and their
empirical confirmations.  The NASWOT row checks an algebraic fact for $x^2$ networks and the others are
empirical agreement, not proofs of exactness.

\begin{table}[h]
\centering
\caption{Theory-empirical connections.}
\label{tab:theory}
\small
\begin{tabular}{llll}
\toprule
\textbf{Result} & \textbf{Prediction} & \textbf{Empirical confirmation} \\
\midrule
Theorem~\ref{thm:det} (cost determinism)
  & Structural model tracks ezkl circuit size
  & exact logrows for $49/73$, within one for $66/73$; $\rho=0.80$ vs.\ ezkl logrows \\
Theorem~\ref{thm:sep} (FLOPs does not determine cost)
  & Ratio $=(\bar f+\kappa_B)/(\bar f+\kappa_A)$
  & ReLU vs.\ $x^2$: $1.00$--$1.06\times$ cells (no separation); unnormalized $x^3{+}x$: same cells, $\logrows$ up $1$--$4$ \\
Proposition~\ref{prop:naswot} (NASWOT degeneracy)
  & $-\infty$ for $x^2$ networks
  & Implementation returns $-\infty$ on all $272$ CIFAR-10 polynomial networks, including the $79$ $x^2$ ones, because it finds no ReLU modules; this does not test the proposition \\
Proposition~\ref{prop:zen} (ZenNAS inflation)
  & Zen $\propto d^L \sigma_x$, anti-correlated
  & ZenNAS $\rho=-0.545$ overall, negative in all 4 families \\
\bottomrule
\end{tabular}
\end{table}


% ============================================================
\section{The zkml-bench Benchmark}
\label{sec:benchmark}
% ============================================================

\subsection{Search Space}

\paragraph{Activation types.}
$\{$ReLU, $x^2$, $x^3{+}x$, $x^2{+}x\}$.  ReLU is the baseline (lookup-based);
$x^2$ is the simplest polynomial; $x^3{+}x$ is a signed
polynomial with better expressivity; $x^2{+}x$ approximates GELU-style
behavior.

\paragraph{Normalization.}
$\{$none, LayerNorm, GroupNorm$\}$.  BatchNorm is excluded: it involves
division by $\sqrt{\mathrm{variance}}$, which is expensive in ZK contrary to
common assumption.

\paragraph{Convolution type.}
$\{$standard conv, depthwise separable (DWS)$\}$.  DWS reduces parameters
and proportionally reduces circuit constraints.

\paragraph{Architecture families.}
MLP, CNN, ResNet-style with skip connections, DepthwiseSepConv-based.

\subsection{Metrics}

Evaluated on three datasets used by NAS-Bench-201~\citep{dong2020nasbench201}:
CIFAR-10 ($n=335$, primary), CIFAR-100 ($n=84$), and ImageNet-16-120
(16$\times$16 downsampled ImageNet, 120 classes).  Metrics are:
\texttt{top1} (30-epoch accuracy), \texttt{constraints} (structural ZK cost
proxy), \texttt{ezkl\_logrows} (actual logrows from ezkl calibration), and
\texttt{ezkl\_mock\_time\_s} (mock proof time in seconds).

\subsection{Zero-Cost Proxies Evaluated}

\begin{table}[h]
\centering
\caption{Zero-cost proxies evaluated in zkml-bench.}
\label{tab:proxies}
\small
\begin{tabular}{lllll}
\toprule
\textbf{Proxy} & \textbf{Reference} & \textbf{Data-free} & \textbf{Poly-compatible} \\
\midrule
Synflow  & \citet{tanaka2020synflow}   & \checkmark & \checkmark \\
NEAR     & \citet{husistein2025near}  & $\times$ (1 batch) & \checkmark \\
ZenNAS   & \citet{lin2021zen}          & \checkmark (Gaussian) & \checkmark \\
ZiCo     & \citet{li2023zico}          & $\times$ (1 batch) & \checkmark \\
MeCo     & \citet{jiang2023meco}       & \checkmark & \checkmark \\
NASWOT   & \citet{mellor2021naswot}    & $\times$ & \textbf{No} (Proposition~\ref{prop:naswot}) \\
\textbf{Structural (ours)} & This work & \checkmark & \checkmark \\
\bottomrule
\end{tabular}
\end{table}

\subsection{Algorithm: Joint Proxy-Guided ZK-ML NAS}

Algorithm~\ref{alg:joint} describes our joint accuracy-cost proxy search,
using Lagrangian scalarization:
\[
J(a) = \mathrm{NEAR}(a) - \lambda \cdot \log\bigl(\ZKcost(a)\bigr),
\]
where $\lambda$ controls the accuracy-cost tradeoff.  Setting $\lambda=0$
recovers pure accuracy prediction; $\lambda\to\infty$ minimizes cost only.

\begin{algorithm}[h]
\caption{Joint Proxy-Guided ZK-ML NAS}
\label{alg:joint}
\begin{algorithmic}[1]
\Require Candidate architectures $\mathcal{A} = \{a_1,\ldots,a_n\}$,
         ZK budget $B$ (max $\logrows$), cost weight $\lambda$
\Ensure Top-$k$ architectures for training
\For{each $a \in \mathcal{A}$}
  \State $C(a) \leftarrow \textsc{StructuralProxy}(a)$
         \hfill $\triangleright$ $O(L \times W)$, no ezkl run, model of Theorem~\ref{thm:det}
  \State $\logrows(a) \leftarrow \lceil\log_2 C(a)\rceil$
  \If{$\logrows(a) > B$}
    \State \textbf{skip} \hfill $\triangleright$ over budget
  \EndIf
  \State Normalize $\mathrm{NEAR}(a)$ and $\log C(a)$ to $[0,1]$
         \hfill $\triangleright$ NEAR ranges $100$--$800$; $\log C$ ranges $13$--$20$
  \State $J(a) \leftarrow \mathrm{NEAR}(a) - \lambda \cdot \log C(a)$
\EndFor
\State \Return architectures sorted by $J(a)$, take top-$k$
\end{algorithmic}
\end{algorithm}

\textbf{Key properties:} The structural proxy in step 2 is a deterministic,
training-free cost model (Theorem~\ref{thm:det}) that matches ezkl-calibrated $\logrows$ exactly for $67\%$ of the
architectures we checked and within one for $90\%$ (Section~\ref{sec:results}).  Normalization in step 6 is critical: without it,
NEAR dominates entirely.  The budget constraint (step 4) is applied to the \emph{estimated}
$\logrows$; the ezkl-calibrated value can be higher (the selected architecture
in Table~\ref{tab:ezkl} is estimated at $21$ and calibrated at $22$).  Total cost:
one model evaluation per candidate, vs.\
an ezkl calibration or proof per candidate.


% ============================================================
\section{What Drives ezkl Circuit Cost}
\label{sec:findings}
% ============================================================

This section reports width-matched measurements of ezkl circuit size.  Skeletons are
identical across columns and only the activation changes.  For each network we run ezkl
(version 9.5.0) \texttt{gen\_settings} and \texttt{calibrate\_settings}
(\texttt{target=resources}) with $40$ random Gaussian calibration inputs and read the
calibrated $\logrows$ and the row count; no proofs are generated.  Weights are random
(seed $0$) and untrained, so calibration ranges reflect untrained networks.  Polynomial
networks are exported without clamping.  Each cell is a single run.

\begin{table}[h]
\centering
\caption{Width-matched ezkl circuit size: total assigned cells (calibrated $\logrows$) for the
  same skeleton with different activations (no normalization, CIFAR-10 input).  Weights and
  calibration data were re-drawn for three seeds; all assignment counts and $\logrows$ are
  identical across the three seeds, so one seed is shown.}
\label{tab:matched}
\small
\begin{tabular}{lrrrr}
\toprule
\textbf{Skeleton} & \textbf{ReLU} & $\mathbf{x^2}$ & $\mathbf{x^2{+}x}$ & $\mathbf{x^3{+}x}$ \\
\midrule
MLP 128-64-32 & 0.407M (19) & 0.408M (19) & 0.408M (19) & 0.404M (23) \\
MLP 256-128 & 0.823M (20) & 0.825M (20) & 0.825M (20) & 0.818M (23) \\
MLP 512-256 & 1.709M (21) & 1.712M (21) & 1.713M (21) & 1.697M (23) \\
CNN 16-32 & 1.734M (21) & 1.832M (21) & 1.857M (21) & 1.879M (23) \\
CNN 16-32-64 & 2.795M (22) & 2.910M (22) & 2.938M (22) & 2.963M (23) \\
\bottomrule
\end{tabular}
\end{table}

\paragraph{ReLU and $x^2$ have the same circuit size at these widths.}
We use ezkl's \texttt{total\_assignments} as the circuit-size measure $C$: for all $45$
runs of ReLU, $x^2$ and $x^2{+}x$ networks (five skeletons, three seeds),
$\lceil\log_2 C\rceil$ equals the calibrated $\logrows$ exactly.  The $x^2$ network has the same calibrated
$\logrows$ as the ReLU network in all $15$ skeleton-seed pairs, and $0.2\%$ more assignments in the MLPs and
$4$--$6\%$ more in the CNNs.  At these widths the
multiply-accumulate rows dominate, and we find no evidence that ReLU is more expensive than
$x^2$.  This contradicts a claim made in an earlier version of this work.  The large cost
differences between architectures in Table~\ref{tab:ezkl} therefore come from layer type
(depthwise and standard convolutions versus MLP layers) and width, not from the ReLU/$x^2$
choice, and the lowest-cost architectures in that table are small polynomial MLPs because
they are small MLPs.

\paragraph{Unnormalized $x^3{+}x$ networks calibrate to $\logrows{=}23$.}
In every unnormalized skeleton, in all three seeds, the $x^3{+}x$ network has about the same number of
assignments as the others ($0.99$--$1.06\times$ ReLU) but is calibrated at $\logrows{=}23$, against $19$--$22$ for the
other activations; ezkl's calibration also selected a lower input scale ($11$ versus $13$).  This
does not hold for every $x^3{+}x$ network: in the $73$-architecture comparison of
Section~\ref{sec:results} the four unnormalized $x^3{+}x$ networks are all at $23$, while the
$13$ with LayerNorm or GroupNorm range from $21$ to $26$ (LayerNorm $22$--$24$, GroupNorm $21$--$26$).  We
did not investigate the mechanism.  The data are consistent with large-range activation outputs forcing
a larger circuit through calibration, but we did not test this.

\paragraph{Runs are not always comparable.}
Re-drawing weights and calibration data did not change any assignment count or $\logrows$ in
Table~\ref{tab:matched}.  Separately, the same architecture can calibrate differently between the
earlier run behind Table~\ref{tab:ezkl} and the fresh run in Table~\ref{tab:proxyezkl}; for example
\texttt{dws\_relu\_none\_16$\times$32$\times$64} is at $\logrows{=}21$ in the earlier run and $20$ in the new
one.  The earlier run's ezkl version, settings and weight state were not recorded, so we cannot say why.

\paragraph{Hypotheses we did not test.}
From reading ezkl's circuit layout code (\texttt{ezkl/src/circuit/ops/layouts.rs}) we expect
that reshape and concatenation add no rows, that the reciprocal square root in LayerNorm is
a costly operation, and that max-pooling costs more than average-pooling.  We did not
measure these and make no claim about them.


% ============================================================
\section{Empirical Results}
\label{sec:results}
% ============================================================

\subsection{Experimental Setup}

\textbf{Architecture families:} MLP (2--4 hidden layers), CNN (2--4 conv
layers), DWS (depthwise separable), ResNet-style.  All with AdaptiveAvgPool +
linear classifier.

\textbf{Training:} 30-epoch SGD ($\mathrm{lr}{=}0.1$ cosine decay,
momentum$=0.9$, weight\_decay$=5{\times}10^{-4}$, batch\_size$=128$).
Polynomial activation models without normalization use $\mathrm{lr}{=}0.01$
with gradient clipping (clip$=1.0$) to prevent explosion.

\textbf{ezkl validation:} Actual ezkl circuit compilation and mock proof times
measured for the architectures of Table~\ref{tab:ezkl} in an earlier run whose ezkl version
and settings were not recorded.  The calibration experiments of Sections~\ref{sec:findings}
and~\ref{sec:results} use ezkl 9.5.0 with \texttt{calibrate\_settings} (\texttt{target=resources})
and $20$--$40$ random calibration inputs.

\textbf{Hardware:} The multi-seed validation (Section~\ref{sec:results}) used one NVIDIA A10G GPU.
The ezkl calibration experiments of Sections~\ref{sec:findings} and~\ref{sec:results} ran on a
24-core CPU workstation.  The hardware for the earlier mock- and real-proof timings was not recorded.

\textbf{Robustness validation:} We additionally retrained 26 targeted CIFAR-10
architectures on one NVIDIA A10G under seeds 42, 43, and 44.  The suite covers
the eight architectures in Table~\ref{tab:ezkl} and activation-matched controls
for ReLU, $x^2$, $x^3{+}x$, and $x^2{+}x$.  Each 30- and 100-epoch run starts
from a fresh initialization and uses its own cosine schedule; the 30th epoch
of a 100-epoch schedule is not reused as a 30-epoch result.  This targeted
experiment is separate from the broader single-seed benchmark sweep.

\subsection{Multi-Seed and Training-Budget Robustness}

Table~\ref{tab:robustness} reports the direct replication of the eight central
accuracy claims.  All values are test top-1 accuracy; multi-seed entries are
mean $\mathbin{\pm}$ sample standard deviation over three seeds.  Every
single-seed value is within $1.2$ percentage points of the corresponding
three-seed 30-epoch mean.  In particular, the joint proxy's $75.0\%$ pick
replicates at $74.73\%\mathbin{\pm}1.61\%$ and reaches
$79.96\%\mathbin{\pm}0.40\%$ after 100 epochs.

The activation-matched controls reinforce that activation effects are
architecture-dependent.  ReLU remains the highest-accuracy choice in each
matched convolutional group, while $x^2{+}x$ is the strongest polynomial
alternative.  In the normalization-free MLP, however, $x^2{+}x$ reaches
$55.67\%$ at 100 epochs versus $53.52\%$ for ReLU.  Pure-$x^2$ DWS controls
remain at chance accuracy under both budgets, so polynomial activations should
not be treated as interchangeable.

\begin{table}[h]
\centering
\caption{Multi-seed replication and longer-training validation of the eight
  architectures used in the accuracy/proving table.  ``Original'' is the
  reported single-seed result; validation columns use fresh training runs.}
\label{tab:robustness}
\small
\begin{tabular}{lrrr}
\toprule
\textbf{Architecture} & \textbf{Original 30 ep.} & \textbf{3-seed 30 ep.} & \textbf{3-seed 100 ep.} \\
\midrule
MLP $x^2{+}x$, none (128/64/32)       & $51.1$ & $50.67\mathbin{\pm}0.52$ & $55.67\mathbin{\pm}0.32$ \\
MLP $x^2{+}x$, LN (128/64/32)         & $52.0$ & $50.87\mathbin{\pm}1.16$ & $54.10\mathbin{\pm}1.77$ \\
DWS ReLU (16/32/64)                   & $62.9$ & $64.10\mathbin{\pm}0.52$ & $67.52\mathbin{\pm}1.12$ \\
DWS ReLU (32/64/128)                  & $69.3$ & $70.48\mathbin{\pm}1.30$ & $75.70\mathbin{\pm}0.71$ \\
DWS ReLU (32/64/128/256)              & $75.0$ & $74.73\mathbin{\pm}1.61$ & $79.96\mathbin{\pm}0.40$ \\
CNN ReLU (16/32/64)                   & $77.1$ & $76.98\mathbin{\pm}0.72$ & $79.80\mathbin{\pm}0.46$ \\
CNN ReLU (16/32/64/128)               & $81.6$ & $81.31\mathbin{\pm}0.53$ & $84.20\mathbin{\pm}0.32$ \\
CNN ReLU (64/128/256)                 & $83.9$ & $83.96\mathbin{\pm}0.16$ & $86.25\mathbin{\pm}0.21$ \\
\bottomrule
\end{tabular}
\end{table}

Across all 26 architecture means, the 30- to 100-epoch Spearman correlation
is $\rho=0.932$.  Across the 77 successful architecture-seed pairs it is
$\rho=0.941$, and the three seed-wise correlations range from $0.929$ to
$0.964$.  Absolute accuracy rises by $5.88$ percentage points on average, so
the longer budget changes accuracy levels more than relative ordering.  The
ordering is strong but not perfect: some deep polynomial CNNs learn late.
Of 156 planned runs, 154 complete; both divergences are the same
$x^3{+}x$ LayerNorm MLP under seed 44, once at each epoch budget, with a
non-finite loss during epoch 1.

\subsection{Proxy Correlations Across Activation Families}

Table~\ref{tab:corr} shows Spearman $\rho$ between each proxy score and top-1
accuracy, stratified by activation type ($n=335$ CIFAR-10 architectures).  The
last two rows report the structural ZK-cost proxy separately against accuracy
and against ZK cost, since these are different targets.

\begin{table}[h]
\centering
\caption{Proxy correlation $\rho$ per activation family; all rows except the
  last are correlations with top-1 accuracy.
  All correlations are Spearman rank correlations.
  N/A (degenerate) indicates our implementation returns $-\infty$ for all architectures
  in that family because it finds no \texttt{nn.ReLU} modules to score (see the remark after
  Proposition~\ref{prop:naswot}).
  ${}^{*}$ $p < 0.05$.  ${}^{\dagger}$ The structural proxy and the sweep's
  ZK cost are the same constraint count, so $\rho{=}1$ is an identity, not a
  validation; the validation of the model is its agreement with ezkl-calibrated
  logrows ($\rho=0.80$, $n=73$; Table~\ref{tab:proxyezkl}).  The structural proxy is a cost model, not an
  accuracy predictor.}
\label{tab:corr}
\small
\begin{tabular}{lrrrr}
\toprule
\textbf{Proxy}
  & \textbf{ReLU} ($n{=}63$)
  & \textbf{$x^2$} ($n{=}79$)
  & \textbf{$x^3{+}x$} ($n{=}102$)
  & \textbf{$x^2{+}x$} ($n{=}91$) \\
\midrule
NEAR        & $+0.828^*$  & $-0.089$   & $+0.026$   & $-0.014$   \\
Synflow     & $+0.079$    & $-0.427^*$ & $+0.120$   & $-0.151$   \\
ZiCo        & $+0.393^*$  & $+0.168$   & $+0.538^*$ & $-0.037$   \\
MeCo        & $+0.416^*$  & $-0.035$   & $+0.034$   & $+0.064$   \\
ZenNAS      & $-0.358^*$  & $-0.162$   & $-0.187$   & $-0.455^*$ \\
NASWOT      & $+0.615^*$  & N/A (degenerate) & N/A (degenerate) & N/A (degenerate) \\
Structural ZK cost vs.\ \emph{accuracy}
            & $+0.650^*$  & $-0.206$   & $-0.246^*$ & $-0.146$   \\
\textbf{Structural ZK cost vs.\ ZK cost}${}^{\dagger}$
            & $+1.000^*$  & $+1.000^*$ & $+1.000^*$ & $+1.000^*$ \\
\bottomrule
\end{tabular}
\end{table}

\textbf{Key findings.}  (1) The structural proxy is not an accuracy predictor ($\rho=+0.65$ for ReLU and
negative or insignificant for the polynomial families); its $\rho=1$ with
sweep ZK cost is an identity.
(2) NASWOT returns $-\infty$ for $272/272$ polynomial CIFAR-10 architectures ($0/63$
for ReLU), consistent with Proposition~\ref{prop:naswot}.  (3) NEAR collapses from
$+0.828$ (ReLU) to $-0.089$ ($x^2$) --- a complete failure on the ZK-ML
design space.  (4) ZenNAS is the most consistently wrong proxy: negative
$\rho$ in all four activation families.  (5) No single proxy works across all
families.

\subsection{Cross-Dataset Validation}

\begin{table}[h]
\centering
\caption{NEAR cross-dataset comparison (same proxy, same code).}
\label{tab:cross}
\small
\begin{tabular}{lrrrr}
\toprule
\textbf{Dataset} & \textbf{NEAR $\rho$ (ReLU)} & \textbf{NEAR $\rho$ ($x^2$)} & \textbf{Gap} \\
\midrule
Our CIFAR-10             & $+0.828^*$ & $-0.089$   & $0.917$ \\
Our CIFAR-100            & N/A ($n{=}3$) & $-0.624^*$ & ---  \\
Our ImageNet-16-120      & $+0.933^*$ & $-0.127$   & $1.060$ \\
\bottomrule
\end{tabular}
\end{table}

The proxy failure appears on both datasets where each family has data.  NASWOT
returns $-\infty$ on every polynomial architecture in our CIFAR-10 sweep ($272$); for $x^2$
networks this is Proposition~\ref{prop:naswot}, and for the others it is because there
are no ReLU patterns to score (see the remark after it).

\subsection{ZenNAS Training Failure Analysis}

\begin{table}[h]
\centering
\caption{ZenNAS score vs.\ training failure rate per architecture family.
  High Zen scores correlate with training instability in BN-free polynomial networks.}
\label{tab:zen}
\small
\begin{tabular}{lrrr}
\toprule
\textbf{Architecture family} & \textbf{Mean ZenNAS} & \textbf{Training failure rate} & \textbf{Mean accuracy} \\
\midrule
ResNet + $x^2$      & $+18.9$ & $57\%$ & $26.5\%$ \\
ResNet + $x^2{+}x$  & $+17.1$ & $56\%$ & $27.9\%$ \\
CNN + $x^2$         & $-11.2$ & $18\%$ & $42.2\%$ \\
CNN + $x^2{+}x$     & $-8.8$  & $9\%$  & $54.2\%$ \\
ReLU (all families) & $-10$ to $-13$ & $0\%$ & $52$--$69\%$ \\
\bottomrule
\end{tabular}
\end{table}

\subsection{Actual ezkl Mock and Real Proof Times}

Table~\ref{tab:ezkl} reports measured ezkl results for key architectures.
\emph{Mock time} is the time of ezkl's mock prover (witness generation plus
constraint check), which does not produce a proof; \emph{prove time} is the time
to generate a real proof and was measured only for the three architectures
shown.  Each value is a single measurement; runs were not repeated and machine load was not
recorded, so times should be read as orders of magnitude, and we do not
claim that time is monotone in $\logrows$ or constraints (for example the
$\logrows{=}21$ DWS network has a mock time of $272$s, similar to the
$\logrows{=}22$ network at $256$s).  The reliable cost measure in this paper is
the calibrated circuit size ($\logrows$ and constraints).

\begin{table}[h]
\centering
\caption{ezkl measurements.  ``est.'' is the structural-proxy estimate of $\logrows$;
  ``cal.\ (earlier)'' is ezkl's calibrated $\logrows$ from an earlier run whose weight
  state was not recorded, with the mock/prove/verify times from that run; ``cal.\ (untrained)''
  is a fresh calibration with random weights (Section~\ref{sec:results}).
  Mock and prove times are single runs; a dash means the quantity was not measured
  (prove and verify times exist only for three architectures).  No mock result was
  recorded for the $\logrows{=}25$ architecture.}
\label{tab:ezkl}
\scriptsize
\begin{tabular}{lrrrrrrr}
\toprule
\textbf{Architecture} & \textbf{top-1} & \textbf{est.} & \textbf{cal.\ (earlier)} & \textbf{cal.\ (untrained)}
  & \textbf{mock (s)} & \textbf{prove (s)} & \textbf{verify (s)} \\
\midrule
mlp\_x2px\_none\_128x64x32 & 51.1\% & 19 & 18 & 19 & 19.1 & 57.5 & 0.20 \\
mlp\_x2px\_layernorm\_128x64x32 & 52.0\% & 19 & 18 & 19 & 19.2 & 46.7 & 0.43 \\
dws\_relu\_none\_16\_32\_64 & 62.9\% & 19 & 21 & 20 & 272.5 & 1339.9 & 4.74 \\
dws\_relu\_none\_32\_64\_128 & 69.3\% & 21 & 22 & 21 & 256.4 & --- & --- \\
dws\_relu\_none\_32\_64\_128\_256 & 75.0\% & 21 & 22 & 22 & 534.4 & --- & --- \\
cnn\_relu\_none\_16\_32\_64 & 77.1\% & 22 & 22 & 22 & 949.5 & --- & --- \\
cnn\_relu\_none\_16\_32\_64\_128 & 81.6\% & 22 & 22 & 22 & 1119.1 & --- & --- \\
cnn\_relu\_none\_64\_128\_256 & 83.9\% & 26 & 25 & 26 & no result & --- & --- \\
\bottomrule
\end{tabular}
\end{table}

The $75.0\%$ architecture (\texttt{dws\_relu\_none\_32$\times$64$\times$128$\times$256}) is estimated at
$\logrows{=}21$ by the structural proxy and calibrated at $22$ by ezkl in both runs.  The
estimate and the calibration disagree for several rows in this table, in both directions, and the
two calibrations of the same architecture also differ for some rows, so the budget in
Table~\ref{tab:search} should be read as a constraint-count budget on the estimate, not as a guarantee
on ezkl logrows.

\textbf{Deployment context.}  We did not measure on mobile hardware, so we make no
claim about on-device feasibility.  The two small polynomial MLPs have
the lowest measured prove times ($47$--$58$s) and the $\logrows{=}21$ DWS ReLU
network took $1340$s, but these are different layer types and sizes, not a controlled activation comparison; a $\logrows{=}25$ circuit has $2^{25}$ rows and
we recorded no proving result for it.

\subsection{Budget-Constrained Architecture Search}

\begin{table}[h]
\centering
\caption{Joint proxy vs.\ individual proxies and random selection in
  budget-constrained search (top-20\% by proxy score).  Budgets are on the
  structural-proxy estimate $\lceil\log_2 C\rceil$, not on ezkl-calibrated
  logrows (the $75.0\%$ architecture is calibrated at $22$, Table~\ref{tab:ezkl}).
  Entries are the best accuracy among selected architectures; the last column's
  time is the single-run mock time.}
\label{tab:search}
\small
\begin{tabular}{lrrrr}
\toprule
\textbf{Strategy}
  & \textbf{est.\ $\logrows{\le}20$}
  & \textbf{est.\ $\logrows{\le}21$}
  & \textbf{est.\ $\logrows{\le}22$}
  & \textbf{Proof time (${\le}22$ winner)} \\
\midrule
NEAR only              & $62.9\%$ & $75.0\%$ & $81.6\%$ & 1119s \\
Synflow only           & $53.1\%$ & $61.0\%$ & $81.6\%$ & 1119s \\
ZK cost only           & $52.0\%$ & $62.9\%$ & $62.9\%$ & 272s  \\
\textbf{Joint proxy (ours)} & $\mathbf{62.9\%}$ & $\mathbf{75.0\%}$ & $\mathbf{75.0\%}$ & $\mathbf{534s}$ \\
Random (avg)           & $55.5\%$ & $62.8\%$ & $73.5\%$ & ---   \\
\bottomrule
\end{tabular}
\end{table}

With an estimated budget of $\logrows{\le}21$, the joint proxy selects the best
architecture in the swept set within that estimated budget ($75.0\%$) while random
selection falls $12.2$ points short ($62.8\%$ on average).  The estimated-budget
winner calibrates to $\logrows{=}22$ in ezkl, so the budget should be read as a
constraint-count budget ($C<2^{21}$), not a guarantee on ezkl logrows.  The joint proxy ties NEAR-only at the two tighter budgets, so its advantage over NEAR
appears only at estimated $\logrows{\le}22$.  There it selects $75.0\%$ accuracy at $534$s
mock time over NEAR's $81.6\%$ at $1119$s: roughly half the single-run mock time
for $6.6$ accuracy points.  Whether that trade is appropriate depends on the
deployment's time budget, which we did not measure.  The per-architecture sweep file behind this
table ($n=335$ trained CIFAR-10 architectures) is not included in the released
repository, so this table cannot be regenerated from the repository alone; the sweep
generator is included, and the ezkl and proxy measurements in
Tables~\ref{tab:ezkl} and~\ref{tab:proxyezkl} are.

\subsection{Structural Proxy Validation}

We compare the structural proxy's estimated $\logrows$ with ezkl's calibrated $\logrows$
on $73$ architectures: a seeded stratified sample of the sweep's MLP and CNN families
(ResNet-style models excluded to keep calibration tractable) plus the architectures
of Table~\ref{tab:ezkl}; all $73$ calibrated without failure.  Networks have random untrained weights (seed $0$) and are calibrated with
random Gaussian inputs, so this tests the cost model, not trained-network behaviour.
The estimate equals ezkl's value exactly for $49$ of $73$ architectures ($67\%$) and is
within one $\logrows$ for $66$ ($90\%$); errors run in both directions (range $-3$ to $+2$).
Table~\ref{tab:proxyezkl} gives the breakdown.  Agreement is strong for ReLU, $x^2$ and
$x^2{+}x$ networks and weak for $x^3{+}x$, where calibration frequently chooses a larger
circuit than the row count implies (Section~\ref{sec:findings}).

\begin{table}[h]
\centering
\caption{Structural-proxy $\logrows$ versus ezkl-calibrated $\logrows$ ($73$ architectures,
  untrained weights).  Spearman $\rho$ with $95\%$ bootstrap CI (2000 resamples).}
\label{tab:proxyezkl}
\small
\begin{tabular}{lrrrr}
\toprule
\textbf{Family} & $n$ & \textbf{exact} & \textbf{within 1} & \textbf{Spearman $\rho$ [95\% CI]} \\
\midrule
All          & 73 & 49 & 66 & $0.80$ $[0.65,0.91]$ \\
ReLU         & 21 & 16 & 21 & $0.92$ $[0.76,0.98]$ \\
$x^2$        & 16 & 12 & 16 & $0.94$ $[0.85,0.99]$ \\
$x^2{+}x$    & 19 & 14 & 17 & $0.93$ $[0.82,0.98]$ \\
$x^3{+}x$    & 17 &  7 & 12 & $0.28$ $[-0.41,0.75]$ \\
\bottomrule
\end{tabular}
\end{table}

Within the CIFAR-10 sweep the structural cost equals the sweep's constraint count by
construction ($\rho=+1.000$ in every family), so that agreement is an identity and not
evidence; Table~\ref{tab:proxyezkl} is the evidence for the model.  The proxy needs no ezkl run;
we did not time it against calibration.


% ============================================================
\section{Related Work}
\label{sec:related}
% ============================================================

\paragraph{ZK-ML provers.}
ezkl~\citep{ezkl2022} is the primary Halo2-based system (our
evaluation target).  Orion/Giza uses Cairo/STARK (no trusted setup, larger
proofs).  RISC Zero and SP1 use RISC-V zkVMs.  None include NAS tooling or
benchmark data.  A recent survey of ZK-based verifiable machine
learning~\citep{peng2025zkmlsurvey} names limited circuit expressiveness and
high proving cost among the main obstacles; architecture-level search for
proving cost is one way to address the latter.

\paragraph{ZK circuit optimization.}
TeleSparse~\citep{maheri2025telesparse} applies post-hoc sparsification and
neural teleportation (to narrow activation ranges and hence lookup tables) to reduce
proof time by $46\%$ at roughly $1\%$ accuracy cost.  This is complementary to
our approach: TeleSparse compresses a fixed ReLU architecture after-the-fact;
we identify ZK-efficient architectures by design before any training.  ZKML
(EuroSys 2024) and ZEN (2021) optimize compiler passes; to our knowledge none
searches over architecture space.  Two NAS works are close in name but not in
goal: PASNet~\citep{peng2023pasnet} searches over polynomial-versus-ReLU activations
for two-party-computation private inference, where the cost model is
communication and latency rather than ZK circuit size; and
zkNAS~\citep{wang2026zknas} uses zero-knowledge proofs to let a server prove
that an outsourced architecture search was run correctly, rather than searching for
architectures that are cheap to prove.

\paragraph{MPC/FHE NAS.}
CryptoNAS~\citep{ghodsi2020cryptonas} minimizes ReLU count for MPC (garbled
circuits).  In MPC, linear layers are free via secret sharing; in ZK-SNARKs,
they are not.  Across our $n{=}335$ sweep, the Spearman correlation between
non-linear operation count (the CryptoNAS proxy) and ZK circuit constraints is
$\rho=0.860$, but Kendall's $\tau=0.693$ --- meaning $16\%$ of architecture
pairs are ranked in opposite order.  CryptoNAS results do not transfer to ZK
deployment.  LinGCN (NeurIPS 2023) linearizes nonlinear layers for homomorphically encrypted
inference, where multiplicative depth is the relevant cost; in our ZK setting the
cost model depends on circuit size across all layers.

\paragraph{Training-free NAS proxies.}
NEAR~\citep{husistein2025near}, AZ-NAS~\citep{lee2024aznas},
ZiCo~\citep{li2023zico}, and MeCo~\citep{jiang2023meco} were developed and
evaluated on standard ReLU-based NAS benchmarks (e.g.\ NAS-Bench-201 and NATS-Bench).
In our benchmark NEAR collapses from $+0.828$ on ReLU to $-0.089$ on $x^2$
architectures, and NASWOT is degenerate on $x^2$ networks (Proposition~\ref{prop:naswot}).
To our knowledge this paper provides the first benchmark evidence of proxy failure in
the ZK-ML design space.

\paragraph{NAS benchmarks.}
NASBench-201~\citep{dong2020nasbench201} is the de facto standard, using
ReLU+BatchNorm exclusively.  NATS-Bench~\citep{dong2021natsbench} likewise
focuses on standard activations.  Neither includes polynomial activations,
BN-free designs, or ZK cost as a metric.  to our knowledge zkml-bench is the first benchmark
to cover these dimensions.


% ============================================================
\section{Limitations and Discussion}
\label{sec:limitations}
% ============================================================

\paragraph{Scale.}
We restrict the benchmark to CIFAR-scale inputs.  Our own extrapolation
(not measured) is that a ResNet-50 on $224{\times}224$ ImageNet would need on the
order of $\logrows{\approx}35$ ($2^{35}\approx 34$ billion rows), which we expect
to be impractical to prove with current hardware; our largest measured circuit is
$\logrows{=}22$, and the most accurate CNN in our sweep needs $\logrows{=}25$.
Results here therefore do not transfer to ImageNet-scale models without new
measurements.

\paragraph{Training epochs.}
We use 30-epoch accuracy as the ground truth.  The literature establishes that
rank correlation with final accuracy plateaus by $\sim$$15$--$20$ epochs for
NAS proxy evaluation~\citep{abdelfattah2021zerocost}.  Our 30-epoch budget is
conservative.  In our targeted 26-architecture validation, independently
trained 30- and 100-epoch architecture means have Spearman $\rho=0.932$
($\rho=0.941$ across 77 paired seed runs; per-seed range $0.929$--$0.964$).
This supports 30 epochs as a high-fidelity ranking signal, while the observed
late learning in some deep polynomial CNNs cautions against treating the
ranking as exact.

\paragraph{Single seed.}
The broad CIFAR-10 ($n=335$) and CIFAR-100 ($n=84$) benchmark sweeps retain
one seed per architecture.  We therefore do not claim that every per-family
proxy coefficient is free of seed noise.  The targeted suite covers 26
architectures with three seeds, including all eight central accuracy/proving
claims: their largest 30-epoch sample standard deviation is $1.61$ percentage
points, and the selected architecture reproduces at
$74.73\%\mathbin{\pm}1.61\%$.  Extending multi-seed coverage to the full
benchmark remains future work.

\paragraph{Practical takeaways for ZK-ML practitioners.}
These are suggestions from our measurements, not tested prescriptions.
(1) Do not assume polynomial activations are cheaper to prove: width-matched, ReLU and
$x^2$ had the same $\logrows$ (Table~\ref{tab:matched}), and unnormalized $x^3{+}x$ networks calibrated to a larger circuit.
The small polynomial MLPs in Table~\ref{tab:ezkl} were cheap because they are small MLPs, and
they have $51$--$52\%$ accuracy against $77$--$84\%$ for the CNNs.
(2) Measure rather than assume: the cost model's estimate differed from ezkl's
calibrated $\logrows$ for $24$ of $73$ architectures (Table~\ref{tab:proxyezkl}).
(3) NEAR, ZenNAS and NASWOT were unreliable on our polynomial architectures; we
did not find a proxy that works across all families.


% ============================================================
\section{Conclusion}
\label{sec:conclusion}
% ============================================================

We showed that deploying neural networks in ZK proof systems creates a
fundamentally different NAS landscape.  The core insight: a ZK circuit cost model is
hardware-independent and available before any training, which makes it a
useful basis for architecture search, with the caveat that the model matches
ezkl-calibrated $\logrows$ exactly for $67\%$ of architectures checked and within one for $90\%$.

Several standard zero-cost proxies were unreliable on the polynomial-activation
architectures we built: NEAR collapses from $\rho{=}{+}0.83$ on our ReLU architectures
to $\rho{=}-0.09$ on $x^2$ architectures; ZenNAS is anti-correlated in all
polynomial families (with a mechanistic explanation in
Proposition~\ref{prop:zen}); NASWOT as implemented returns $-\infty$ on all $272$ polynomial CIFAR-10 networks because it finds no
ReLU modules (Proposition~\ref{prop:naswot} proves degeneracy for $x^2$ only), and ZiCo is the one proxy
that partly worked ($\rho=+0.54$ on $x^3{+}x$, but $-0.04$ on $x^2{+}x$).

Our width-matched ezkl measurements give a negative result for the common assumption
that polynomial activations are cheaper to prove: ReLU and $x^2$ networks have the same
calibrated $\logrows$ at identical widths, and $x^3{+}x$ networks calibrate to a larger
circuit.  In our measurements the selected DWS+ReLU architecture has a mock-proof time of
$534$s versus $1119$s for the largest feasible CNN ($75.0\%$ vs.\ $81.6\%$ accuracy), and
avoids the $\logrows{=}25$ circuit of the most accurate CNN, for which we recorded no result.

\paragraph{Limitations of the new measurements.}
The proxy-versus-ezkl and width-matched experiments use untrained random weights, random
Gaussian calibration inputs, one ezkl version (9.5.0), one seed and single runs, and they
cover MLPs and small CNNs only; calibration is known to depend on weights and data
(Section~\ref{sec:findings}).  They do not cover transformers, residual networks, or trained
networks, and they measure circuit size, not proving time.

The zkml-bench code, the sweep generator, the ezkl measurement scripts, the result
files behind Tables~\ref{tab:ezkl} and~\ref{tab:proxyezkl}, and the multi-seed validation
summary are at \texttt{https://github.com/npow/zk-circuit-aware-nas} (tag
\texttt{v1}).  The per-architecture sweep results behind the correlation and
budget-search tables, and the raw multi-seed job output, are not in that repository.


% ============================================================
\begin{thebibliography}{99}
% ============================================================

\bibitem[Abdelfattah et al.(2021)]{abdelfattah2021zerocost}
Abdelfattah, M.~S., Mehrotra, A., Dudziak, \L., and Lane, N.~D.
\newblock Zero-cost proxies for lightweight nas.
\newblock In \emph{International Conference on Learning Representations}, 2021.

\bibitem[Dong \& Yang(2020)]{dong2020nasbench201}
Dong, X. and Yang, Y.
\newblock {NAS-Bench-201}: Extending the scope of reproducible neural
  architecture search.
\newblock In \emph{International Conference on Learning Representations}, 2020.

\bibitem[Dong et al.(2021)]{dong2021natsbench}
Dong, X., Liu, L., Musial, K., and Gabrys, B.
\newblock {NATS-Bench}: Benchmarking {NAS} algorithms for architecture topology
  and size.
\newblock \emph{IEEE Transactions on Pattern Analysis and Machine Intelligence},
  2021.

\bibitem[Ghodsi et al.(2020)]{ghodsi2020cryptonas}
Ghodsi, Z., Veldanda, A.~K., Reagen, B., and Garg, S.
\newblock {CryptoNAS}: Private inference on a relu budget.
\newblock In \emph{Advances in Neural Information Processing Systems}, 2020.

\bibitem[Husistein et al.(2025)]{husistein2025near}
Husistein, R.~T., Reiher, M., and Eckhoff, M.
\newblock {NEAR}: A training-free pre-estimator of machine learning model
  performance.
\newblock In \emph{International Conference on Learning Representations}, 2025.

\bibitem[Jiang et al.(2023)]{jiang2023meco}
Jiang, T., Wang, H., and Bie, R.
\newblock {MeCo}: Zero-shot {NAS} with one data and single forward pass via
  minimum eigenvalue of correlation.
\newblock In \emph{Advances in Neural Information Processing Systems}, 2023.

\bibitem[Lee \& Ham(2024)]{lee2024aznas}
Lee, J. and Ham, B.
\newblock {AZ-NAS}: Assembling zero-cost proxies for network architecture
  search.
\newblock In \emph{IEEE/CVF Conference on Computer Vision and Pattern
  Recognition}, 2024.

\bibitem[Li et al.(2023)]{li2023zico}
Li, G., Yang, Y., Bhardwaj, K., and Marculescu, R.
\newblock {ZiCo}: Zero-shot nas via inverse coefficient of variation on
  gradients.
\newblock In \emph{International Conference on Learning Representations}, 2023.

\bibitem[Lin et al.(2021)]{lin2021zen}
Lin, M., Wang, P., Sun, Z., Chen, H., Sun, X., Qian, Q., Li, H., and Jin, R.
\newblock {Zen-NAS}: A zero-shot nas for high-performance image recognition.
\newblock In \emph{International Conference on Computer Vision}, 2021.

\bibitem[Mellor et al.(2021)]{mellor2021naswot}
Mellor, J., Turner, J., Storkey, A., and Crowley, E.~J.
\newblock Neural architecture search without training.
\newblock In \emph{International Conference on Machine Learning}, 2021.

\bibitem[Maheri et al.(2025)]{maheri2025telesparse}
Maheri, M.~M., Haddadi, H., and Davidson, A.
\newblock {TeleSparse}: Practical privacy-preserving verification of deep
  neural networks.
\newblock \emph{Proceedings on Privacy Enhancing Technologies}, 2025(4):861--880;
  arXiv:2504.19274.

\bibitem[Peng et al.(2023)]{peng2023pasnet}
Peng, H., Zhou, S., Luo, Y., Xu, N., Duan, S., Ran, R., Zhao, J., Wang, C., Geng,
  T., Wen, W., Xu, X., and Ding, C.
\newblock {PASNet}: Polynomial architecture search framework for two-party
  computation-based secure neural network deployment.
\newblock In \emph{ACM/IEEE Design Automation Conference (DAC)}, 2023;
  arXiv:2306.15513.

\bibitem[Peng et al.(2025)]{peng2025zkmlsurvey}
Peng, Z., and others.
\newblock A survey of zero-knowledge proof based verifiable machine learning.
\newblock \emph{arXiv preprint arXiv:2502.18535}, 2025.

\bibitem[Tanaka et al.(2020)]{tanaka2020synflow}
Tanaka, H., Kunin, D., Yamins, D.~L.~K., and Ganguli, S.
\newblock Pruning neural networks without any data by iteratively conserving
  synaptic flow.
\newblock In \emph{Advances in Neural Information Processing Systems}, 2020.

\bibitem[Wang et al.(2026)]{wang2026zknas}
Wang, H., Zhang, D., Jiang, T., Wang, D., Guo, F., Wang, J., and Guo, Y.
\newblock {zkNAS}: Secure and efficient outsourced-{NAS} with zero-cost proxies.
\newblock Springer Nature Singapore, 2026.
  doi:10.1007/978-981-92-0375-8\_31.

\bibitem[zkonduit(2022)]{ezkl2022}
zkonduit.
\newblock {ezkl}: zero-knowledge machine learning inference toolchain
  (software).
\newblock \url{https://github.com/zkonduit/ezkl}.

\end{thebibliography}

\end{document}
